\section{Discussion}

\subsection{No Link Between Drug Resistance \& Immunogenicity}

Currently, It is well established that one mechanism for which \ac{hiv1} evades the immune system are mutations that inhibit the body's natural ability to display antigenic peptides via \ac{hla1} \cite{varela-rohenaControlHIV1Immune2008}. Furthermore, it is also well established that \ac{hiv1} is, more and more, developing mutations capable of reducing the effectiveness of a variety of drugs \cite{clutterHIV1DrugResistance2016}. Lastly,

It then may be interesting to see if there does exist a general trend of any sort with regards to \acp{drm} in \ac{hiv1} epitope regions. From figure \ref{fig:overall_hist_axs}, we see that the distribution over \ac{pelrank} has proportionately, shifted towards higher rankings when compared to the \acp{naepi}, and thus a lower likelihood of the epitopes successfully binding with \ac{hla1} resulting in the suppression of viral peptide presentation, while not disturbing the functionality of the \ac{hla1} in its natural state. However, these results assume that all hypothetical epitopes are virally viable. In this scenario, \ac{hiv1} would not only be resistant towards the majority of known \ac{ard}, but also be very accomplished in evading the host's immune response thus producing monster \ac{hiv1}. However, it has been established that both forms of mutation, either to escape an effective immune response or to resist the effects of \acp{ard}, will typically result in some form of penalty to the fitness of the \ac{hiv1} strand \cite{clutterHIV1DrugResistance2016,woodmanHIVMolecularEpidemiology2009,leslieHIVEvolutionCTL2004}. From this then, for a host undergoing \ac{art} treatment, the \ac{hiv1} may be said to be undergoing selective pressures from the host immune system, the drug, and its own functional viability, that is to say, the incurring of too many mutations may be deleterious to the innate functionality of proteins encoded for by the strand of \ac{hiv1} itself. It is this last broad selective pressure that is likely preventing such a monster \ac{hiv1} from actually forming. In other words, the fitness burden of maintaining such a degree of \acp{drm} within the \ac{hiv1} genome must be deleterious to its functionality.

Thus, a more realistic distribution would involve filtering the generated epitope set solely for generated epitopes with the same, or less \acp{drm} than that of the greatest number of \acp{drm} in the set of \acp{naepi}. By doing this, are essentially naively assuming for every \ac{drm}, the \ac{drm} can be said to carry 1 fitness burden unit. Then, by resolving for the maximum number of \acp{drm} in the set of \acp{naepi}, we essentially attempt to find the upper viable limit of \ac{drm} based fitness burden units per epitope where going over is less optimal then being more susceptible due to not having a given \ac{drm}. The resulting figure \ref{fig:overall_hist_axs_capped} indicates that in the case that such an assumption holds relatively true, the overall effect of adding at most two \acp{drm} to any given epitope does not affect greatly affect the population of epitope's capacity to bind or rather, avoid the activation of \ac{hla1}, and thus, does not affect \ac{hiv1}'s immunological escape capacities via the suppressing of peptide presentation via \ac{hla1}.

\subsection{Potential Increased Immunogenicity with HLA-B*15:01}

Despite the overall lack of association between the presence of at most 2 \acp{drm} in epitopes, by stratifying the data and representing the data in the form of a \ac{kde}, we can firstly, see the initial conclusion of the overall distribution of \ac{pelrank} not greatly affecting remains visible in each subplot of figure \ref{fig:pairplot_binding_mut_orig}. I.e., the figure plots the \ac{pelrank} of all generated epitopes and the \ac{pelrank} of the \ac{naepi} the generated epitope originated from (root epitope). This, when not accounting for the stratification of different \ac{hla1} alleles, we see that as the rank of root epitopes increases over the horizontal axis, the rank of the mutated epitopes grows at a similar rate with some deviation. However, stratifications that result in distributions that do not follow this trend also exist. Notably, there does seem to exist some sets of root epitopes where the introduction of \acp{drm} resulted in a decrease of \ac{pelrank} thus increasing binding strength to the respective allele. Specifically, of the \ac{insti} class, we see this in the HLA-B*15:01 allele. In contrast, and potentially as expected, the majority of the root epitopes, when introduced with \acp{drm}, result in an increase in the spread of distribution towards an increased \ac{pelrank}, thus, indicating that many \acp{drm} may up-regulate immune escape mechanisms assuming that the combination of \acp{drm} are viable.

HLA-B*15:01 is not prevalently discussed with regards to \ac{hiv1}. Interestingly, one paper has previously associated HLA-B*15:01 with an increased susceptibility to \ac{hiv1} in a Turkish population \cite{darbasRelationshipHLABAlleles2022}. However, this study is not regarding \ac{hiv1} susceptibility, but rather, the effect of \acp{drm} on the immunogenicity of HIV-1. In other words, \ac{hiv1}, in acquiring \acp{drm} to counteract \ac{insti} drugs, may be more susceptible to host HLA-B*15:01 host immune response. Alternatively, if it is true that HLA-B*15:01 increases susceptibility to HIV-1 patients, it may also be true that \ac{insti} drugs will be more effective for those patients in that \ac{insti} drugs may place selective pressure on the \ac{hiv1} strand to develop drug resistance towards \ac{insti}, however, in doing so, may increase such a hypothetical strand's \ac{hiv1} susceptibility to HLA-B*15:01 based host immune response.

\subsection{Specific Drug Resistance Mutations Affects on Immunogenicity}

To begin, we will look at the effects of HLA-B*07:02 on \ac{insti} class \acp{ard}. Notably, multiple papers discuss the protective properties of HLA-B*07:02 towards \ac{hiv1} \cite{darbasRelationshipHLABAlleles2022,sundaramurthiHLAB2705specific2013}. Our results indicate several combinations of \acp{drm} that may be associated with increased immunogenicity by increasing the \ac{pelrank}, and thus undermining potential previously \ac{hla1} mediated immune escapes. The association with a strong increase in immunogenicity is corroborated by the \ac{hivdrmdb} where it indicates such a mutation is rather rare. This rareness may be justified due to its high penalty to \ac{hiv1}'s immune escape capabilities. Furthermore, all mutations within this general region are categorized as highly resistant to most \ac{insti} drugs (table \ref{table:significantly_associated_mutation_drm_scores} for resistance scores and table \ref{table:hivdrmdb_categorizations} for interpreting the scores).

Continuing, HLA-B*15:01 appears again where specifically, mutations set E92Q and Q95K, appear together, and E92Q itself, is associated with a relatively dramatic increase in immunogenicity with high drug resistance capabilities. A separate paper has replicated this result, concluding similarly, however, for a specific peptide \cite{rahmanRaltegravirElvitegravirresistanceMutation2014}. Furthermore, of note, HLA-B*15:01, which was discussed previously as part of the analysis of the joint \ac{kde} plot, supports each other in that both seem to indicate an increase in \ac{hla1} \ac{pelrank} decrease, and thus, a binding strength increase, ultimately, predicting an increase of immunogenicity.

Lastly, N88T and V82M protease epitope mutations were found to be statistically significantly associated with an increase in predicted immunogenicity towards the HLA-B*07:02 \ac{hla1}. Interestingly, \ac{hivdrmdb} notes that N88T is considered extremely rare while only conferring minor resistance towards select \ac{pi} drugs. Similarly, V82M is considered uncommon by \ac{hivdrmdb}. The rareness corroborates our results as we indicate, as a result of the two mutations, a strong increase in immunogenicity. Furthermore, the combination of the two only significantly confers strong resistance towards Nelfinavir according to \ac{hivdrmdb}.
